\chapter{Conclusioni}

Questa tesi ha affrontato il problema del superamento sicuro di un ostacolo nella guida autonoma, sviluppando un framework decisionale integrato in CARLA Leaderboard e confrontando due approcci di controllo diversi, PID e MPC, a parità di scenari, sensori e logica di alto livello. Il lavoro ha mostrato che una macchina a stati relativamente semplice, se supportata da soglie dinamiche di sicurezza, analisi della corsia laterale e override locali di emergenza, è sufficiente per governare in modo coerente l'intera sequenza della manovra: rilevamento del blocco, attesa, cambio corsia, superamento e rientro.

Dal punto di vista implementativo, il contributo principale del lavoro consiste nell'aver trasformato componenti eterogenei e inizialmente separati in un sistema sperimentale completo, capace di produrre un comportamento concreto e misurabile nello scenario di sorpasso. Sono stati infatti definiti scenari rappresentativi, sviluppata una macchina a stati dedicata, integrati sensori e regole di sicurezza, adattati i controller al caso di studio, costruita una pipeline di logging frame-by-frame e introdotta una procedura di esportazione automatica delle metriche. Senza questo lavoro di integrazione, taratura ed estensione, i moduli di partenza non sarebbero stati sufficienti a fornire un risultato utile né un confronto sperimentale difendibile.

I risultati sperimentali mostrano che, negli scenari considerati, il controllore PID è risultato mediamente più preciso e più rapido nel completare la manovra. L'errore laterale medio è stato inferiore rispetto all'MPC, così come il tempo complessivo di manovra e la distanza minima dall'ostacolo si sono mantenuti su valori competitivi o migliori. L'MPC ha mostrato invece un vantaggio qualitativo importante: il comando di sterzo è risultato più regolare e la traiettoria più fluida, a conferma del fatto che l'approccio predittivo riesce a controllare meglio la variazione del comando. Questo beneficio, tuttavia, è stato ottenuto al prezzo di un costo computazionale nettamente superiore.

Nel contesto specifico di questa tesi, quindi, l'MPC non è risultato migliore del PID in senso assoluto, ma diverso. Il PID si è dimostrato una soluzione molto efficace, semplice e leggera per gli scenari studiati, mentre l'MPC ha evidenziato un potenziale maggiore soprattutto sul piano della fluidità e della gestione esplicita dei vincoli. Di conseguenza, il valore dell'MPC appare particolarmente interessante in prospettiva futura, cioè in scenari più complessi, con dinamiche più ricche, vincoli più severi e necessità più forti di anticipazione. Il lavoro svolto fornisce dunque una base concreta per ulteriori estensioni, sia sul lato della percezione e della pianificazione sia sul lato del controllo avanzato.
